\documentclass[12pt]{article}
\usepackage[margin=1in]{geometry}
\usepackage[T1]{fontenc}
\usepackage[utf8]{inputenc}
\usepackage[english]{babel}
\usepackage{graphicx}
\usepackage{subcaption}
\usepackage{float}
\usepackage[section]{placeins}
\renewcommand{\figurename}{Figure}
% \usepackage[turkish]{babel}
\usepackage[expansion=false,protrusion=true]{microtype}
\usepackage{newtxtext}
\usepackage{newtxmath}
\let\Bbbk\relax % newtxmath ile amssymb çakışmasını engelle
\usepackage{amsmath,amssymb}
\usepackage[hidelinks]{hyperref}

\title{ARTPS: Depth-Enhanced Hybrid Anomaly Detection and Learnable Curiosity Score for Autonomous Rover Target Prioritization}
\author{Poyraz BAYDEMİR}
\date{}

\begin{document}
\maketitle

\tableofcontents
\newpage

\begin{abstract}
This work presents an integrated approach for \emph{on-site} prioritization of scientifically interesting targets in autonomous exploration scenarios on planetary surfaces. The proposed ARTPS architecture consists of single-image depth estimation, multi-component anomaly fusion of image and depth cues, engineering-focused target localization, and learnable \emph{curiosity} score components. The system increases sensitivity to small and nearby objects while maintaining stable performance in distant field details; it addresses shadow, specularity, and low texture effects that challenge field operations through field-specific signals combined with morphological and photometric cues.

From a software engineering perspective, the application has an execution plan that considers edge computing constraints, memory/energy budgets, timing requirements, and communication latency together. The image processing pipeline and depth enhancement steps are parametric and can be \emph{fine-tuned} for field conditions through the user interface. System outputs are designed to be explainable to support operator decisions: numbered regions on the combined anomaly map are mapped one-to-one with criteria presented in the right diagnostic panel, and selected targets are visually highlighted.

From a scientific perspective, the \emph{curiosity} score is defined as a normalized combination of known value, reconstruction difference, combined anomaly density, depth variance, and roughness, and is calibrated for field conditions through regularized learning. In evaluation; AUROC/AUPRC for anomaly discrimination, relative error and accuracy thresholds for depth, nDCG and ranking correlations for ranking quality are used. Results show that small-near object sensitivity increases, distant field details are preserved, and shadow/specular suppression significantly reduces false alarms. The work is presented with methodological details, experimental setup, ablations, and reproducibility principles.
\end{abstract}

\section{Giriş}
Gezegen keşif görevlerinde hedef seçimi ve sıra planlama, veri toplama bütçeleri ile bilimsel değer arasındaki dengeyi belirler. Yeryüzü ile iletişimin gecikmeli ve kesintili olması, keşif platformlarının (ör. robotik gezgin) sahada daha fazla özerklik üstlenmesini gerektirir. Bu bağlamda, \textbf{hedef önceliklendirme} ve \textbf{görev zamanlaması} süreçlerinin, çevreden türetilen algısal ipuçlarıyla beslendiği, baştan sona tutarlı bir \emph{saha zekâsı} mimarisine ihtiyaç vardır. Operasyonel kısıtlar; bant genişliği sınırlı telemetri, veri gecikmesi, enerji kısıtları, termal koşullar ve bellek/işlem bütçeleri şeklinde kendini gösterir.

Saha koşullarında ilginç hedeflerin tanımlanması salt \emph{anomali} sinyali ile sınırlı değildir. Bir hedef, jeolojik bağlamda anlamlı bir düzene işaret edebilir (ör. katmanlaşma, tortulanma), ayrıca hedefin çevresel ilişkileri (yerel eğim, z-normal yakınsaması, planarlık/pürüzlülük dağılımı) göz önüne alındığında bilinen değer taşıyabilir. Bu nedenle, \emph{bilinen değer} ve \emph{yenilik/anomali} sinyallerinin dengeli bir bileşimi, bilimsel getiri açısından daha tutarlı bir strateji sunar.

Bu makalede sunulan ARTPS, tek görüntüden derinlik kestirimi ile görüntü/derinlikten türetilen çoklu anomali göstergelerini birleştirerek \emph{ilginçlik} (curiosity) odaklı bir sıralama üretir. Tasarım hedefleri: (i) küçük ve yakın nesnelerde yüksek duyarlılık, (ii) uzak bölgelerde ayrıntıların korunması, (iii) gölge/parlama ve düşük doku gibi sahaya özgü yan etkilerin bastırılması, (iv) sonuçların mühendislik bakış açısından açıklanabilir görseller ve tanı panelleriyle sunulması ve (v) uca yakın hesap kısıtlarını gözeten verimli bir gerçekleştirimdir. Yazılım mimarisi, modüler tasarıma ve \emph{uçta çalıştırılabilirlik} ilkesine uygun olarak inşa edilmiştir.

\section{İlgili Çalışmaların Kısa Özeti}
\textbf{Tek görüntü derinlik kestirimi:} Son yıllarda yoğun kestirim için dönüştürücü tabanlı mimariler ve bunların farklı çözünürlüklerde ince ayarlanmış sürümleri yaygınlaşmıştır. Kenar rehberli iyileştirme, hızlı küresel düzgünleştirme ve ağırlıklı medyan gibi sonradan işleme adımlarıyla geometrik tutarlılık ve ince ayrıntılar güçlendirilebilir. Tek görüntü derinlik yöntemlerinin sahaya aktarımı sırasında, fotometrik tutarsızlıklar (pozlama/dinamik aralık), perspektif bozulmalar ve çok ölçekli doku kayıpları dikkatle yönetilmelidir.

\textbf{Endüstriyel görsel anomali:} Bir uçta yeniden-oluşturma temelli (sıkıştırıcı-açıcı mimariler), diğer uçta öz nitelik istatistiği temelli (yamaya ait çok değişkenli dağılımın modellenmesi) yaklaşımlar ile anı bankası ve komşuluk ilişkilerine dayalı (öz uzayında k-en yakın komşu) yöntemler bulunmaktadır. Çok ölçekli gradyan, laplasyen ve DoG benzeri filtrelerle ayrıntı bileşenleri zenginleştirilirken, düşük varyanslı alanlar ve spekülarite/gölge kaynaklı sahte ipuçları bastırılır.

\textbf{Çoklu ipucu füzyonu:} Görüntü ve derinlikten gelen ipuçlarının normalize edilip ağırlıklı toplanması; gölge/speküler bastırma ve histerezis eşiğiyle gürültü kontrolü yaygın stratejilerdir. Yerelleştirme tarafında döndürülmüş dikdörtgenler, IoU tabanlı bastırma ve kutu birleştirme pratik çözümler sunar. Füzyonun başarısı, bileşenler arası ölçek uyumu ve alan bağlamını gözeten ağırlıklandırmaya bağlıdır.

\textbf{Curiosity odaklı keşif:} Keşif kazanımını artırmak için \emph{bilinen değer} ile \emph{anomali/yenilik} sinyallerini dengeleyen, öğrenilebilir ağırlıklara sahip sıralayıcılar önerilmiştir. Saha koşullarına göre veri güdümlü ayar yapılması öne çıkar; ayrıca açıklanabilirlik ve belirsizlik boyutları, operasyonel karar süreçlerinde güvenliği artırır.

\section{Yöntem}
ARTPS boru hattı, hibrit sistem yaklaşımı ile tasarlanmış olup, farklı veri kaynaklarından gelen bilgileri entegre ederek daha güvenilir sonuçlar üretir. Bütün bileşenler uçta çalıştırılabilir şekilde modülerleştirilmiş; bellek ayak izi, işlem süresi ve sayısal kararlılık açısından profil edilmiştir.

Hibrit sistem yaklaşımının temel prensibi, farklı veri kaynaklarından gelen bilgilerin birleştirilmesi ile tek bir yöntemin sınırlamalarını aşmaktır. Bu yaklaşım, görüntü tabanlı anomali tespiti, derinlik kestirimi ve diğer sensör verilerini entegre ederek daha güvenilir sonuçlar üretir.

Hibrit sistem çalışma mantığı şu şekildedir: (1) \emph{Veri Toplama Aşaması}: Farklı sensörlerden (görüntü, derinlik, spektral) veri toplanır, (2) \emph{Ön İşleme Aşaması}: Her veri kaynağı kendi özelliklerine göre optimize edilir, (3) \emph{Füzyon Aşaması}: Ağırlıklı toplam ile bilgiler birleştirilir, (4) \emph{Karar Aşaması}: Birleştirilmiş bilgiler kullanılarak final karar verilir. Bu yaklaşım, tek bir yöntemin zayıf yanlarını diğer yöntemlerin güçlü yanları ile telafi eder.

\subsection{Sistem Mimarisi ve Yazılım Tasarımı}
ARTPS, üç ana katmandan oluşan modüler bir mimariye sahiptir: (i) \emph{veri işleme katmanı} (görüntü iyileştirme, derinlik kestirimi), (ii) \emph{analiz katmanı} (anomali tespiti, füzyon) ve (iii) \emph{karar katmanı} (curiosity skoru, yerelleştirme). Her katman, bağımsız olarak test edilebilir ve performans profili çıkarılabilir şekilde tasarlanmıştır.

\begin{table}[h]
\centering
\begin{tabular}{|l|l|l|l|}
\hline
\textbf{Katman} & \textbf{Bileşenler} & \textbf{Teknoloji} & \textbf{Çıktı} \\
\hline
Veri İşleme & Görüntü İyileştirme & OpenCV, PIL & Enhanced Image \\
& Derinlik Kestirimi & PyTorch, ViT & Depth Map \\
\hline
Analiz & Anomali Tespiti & PaDiM, Autoencoder & Anomaly Maps \\
& Füzyon & Weighted Fusion & Combined Map \\
\hline
Karar & Merak Skoru & MLP, Ridge Regression & Priority Score \\
& Yerelleştirme & OpenCV, NMS & Bounding Boxes \\
\hline
\end{tabular}
\caption{ARTPS Sistem Mimarisi}
\end{table}

Yazılım mimarisi, Python tabanlı PyTorch ekosistemi üzerine inşa edilmiştir. Modüler tasarım, her bileşenin ayrı ayrı optimize edilmesine ve sahaya özgü konfigürasyonlara olanak tanır. Bellek yönetimi, dinamik bellek tahsisi ve gradyan hesaplama optimizasyonları ile uç cihazlarda verimli çalışma sağlanır.



\subsection{Girdi iyileştirme}
Görsel kaliteyi artırmak ve anomali/kenar ipuçlarını belirginleştirmek amacıyla otomatik bir zincir uygulanır: çözünürlük eşitleme, gürültü azaltma, histogram sınırlamalı yerel kontrast artırma, gama dengeleme ve hafif keskinleştirme. Sınır koruyucu filtreler (kenar rehberli/ortam farkındalıklı) ile aşırı düzleştirme engellenir. Bu adımlar, derinlik kestiriminde kenar tutarlılığını ve anomali bileşenlerinin ayrışmasını iyileştirir. Fotometrik olarak, gölge/speküler izleri belirlemek üzere aydınlık ve doygunluk kanallarında eşik/variyans tabanlı göstergeler hesaplanır.

Görüntü iyileştirme zinciri, aşağıdaki matematiksel işlemleri sıralı olarak uygular:

\begin{enumerate}
\item \textbf{Çözünürlük eşitleme:} Giriş görüntüsü \(I \in \mathbb{R}^{H \times W \times 3}\) için, hedef çözünürlük \(H' \times W'\)'ye bicubic interpolasyon ile yeniden boyutlandırma:
\[ I' = \text{resize}(I, (H', W'), \text{method='bicubic'}) \]

\item \textbf{Gürültü azaltma:} Bilateral filtre ile kenar koruyucu yumuşatma:
\[ I_{\text{denoised}}(p) = \frac{1}{W_p} \sum_{q \in \mathcal{N}(p)} I(q) \cdot w_s(p,q) \cdot w_r(I(p), I(q)) \]
Burada \(w_s\) uzamsal ağırlık, \(w_r\) renk ağırlığı ve \(W_p\) normalizasyon faktörüdür.

\item \textbf{Histogram sınırlamalı kontrast artırma:} Yerel histogram eşitleme ile kontrast iyileştirme:
\[ I_{\text{enhanced}}(p) = \text{CLAHE}(I_{\text{denoised}}(p), \text{clip\_limit}=2.0, \text{tile\_grid\_size}=(8,8)) \]

\item \textbf{Gama dengeleme:} Adaptif gama düzeltme:
\[ \gamma = \frac{\log(0.5)}{\log(\mu_I + \epsilon)}, \quad I_{\text{gamma}}(p) = I_{\text{enhanced}}(p)^{\gamma} \]
Burada \(\mu_I\) görüntünün ortalama yoğunluğu ve \(\epsilon\) sayısal kararlılık için küçük bir değerdir.
\end{enumerate}

\begin{figure}[H]
  \centering
  \begin{subfigure}[t]{0.49\textwidth}
    \centering
    \includegraphics[width=\linewidth]{figures/non_enhanced.jpg}
    \caption{Sisli/ham görüntü}
  \end{subfigure}\hfill
  \begin{subfigure}[t]{0.49\textwidth}
    \centering
    \includegraphics[width=\linewidth]{figures/enhanced.jpg}
    \caption{Sis giderme + iyileştirme sonrası}
  \end{subfigure}
  \caption{Girdi iyileştirme: sis giderme ve fotometrik düzenleme sonrası rover ve yüzey ayrımının netleşmesi.}
  \label{fig:dehaze-enhancement}
\end{figure}
\FloatBarrier

\subsection{Tek görüntü derinlik kestirimi}
Yoğun derinlik haritası, dönüştürücü tabanlı tek-görüntü derinlik mimarileri sınıfında yer alan bir modelle elde edilir. İnce ayrıntıların korunması için: (i) kenar rehberli kılavuz filtreleme, (ii) hızlı küresel düzgünleştirme ve (iii) ağırlıklı medyan filtre seçmeli olarak uygulanır; (iv) test zamanı artımı kapsamında yatay çevirimler füzyonla birleştirilir. Derinlikten türetilen özellikler; istatistikler (ortalama, varyans, çarpıklık, basıklık), eğim büyüklüğü, yüzey normali yakınsaması, planarlık ve pürüzlülük gibi jeomorfolojik özetleri kapsar. Uzak alanlarda ayrıntıların korunması için, derinlik gradyanı tabanlı \emph{zayıf yapılandırılmış} bölgelerde histerezis benzeri birleşik eşikleme uygulanır.

Derinlik kestirimi için kullanılan dönüştürücü mimarisi, Vision Transformer (ViT) tabanlı bir encoder-decoder yapısıdır. Giriş görüntüsü \(I\) için, patch embedding ve pozisyon kodlaması ile başlayan süreç:

\[ E = \text{PatchEmbed}(I) + \text{PositionalEncoding} \]

Patch embedding işlemi:
\[ \text{PatchEmbed}(I) = \text{Linear}(\text{Reshape}(I, (H \cdot W, P^2 \cdot C))) \]

Burada \(P\) patch boyutu, \(C\) kanal sayısı, \(H\) ve \(W\) görüntü boyutlarıdır.

Multi-head self-attention mekanizması ile özellik çıkarımı:
\[ \text{Attention}(Q,K,V) = \text{softmax}\left(\frac{QK^T}{\sqrt{d_k}}\right)V \]

Multi-head attention için:
\[ \text{MultiHead}(Q,K,V) = \text{Concat}(\text{head}_1, \ldots, \text{head}_h)W^O \]

Burada her head:
\[ \text{head}_i = \text{Attention}(QW_i^Q, KW_i^K, VW_i^V) \]

Derinlik haritası \(D \in \mathbb{R}^{H \times W}\) için, decoder katmanları:
\[ D = \text{Decoder}(\text{Encoder}(E)) \]

Feed-forward network:
\[ \text{FFN}(x) = W_2 \cdot \text{ReLU}(W_1 \cdot x + b_1) + b_2 \]

Derinlik iyileştirme adımları:
\begin{itemize}
\item \textbf{Kenar rehberli filtreleme:} Görüntü kenarları \(E_I\) kullanılarak derinlik kenarları \(E_D\) iyileştirilir:
\[ E_D'(p) = E_D(p) \cdot \exp(-\alpha \|E_I(p)\|_2) \]

\item \textbf{Küresel düzgünleştirme:} Poisson denklemi çözümü ile tutarlılık:
\[ \nabla^2 D' = \nabla \cdot \nabla D \]

\item \textbf{Ağırlıklı medyan filtre:} Uzamsal ağırlıklarla kenar koruyucu yumuşatma:
\[ D_{\text{filtered}}(p) = \text{median}\{D(q) \cdot w(p,q) : q \in \mathcal{N}(p)\} \]
\end{itemize}

\begin{figure}[H]
  \centering
  \begin{subfigure}[t]{0.7\textwidth}
    \centering
    \includegraphics[width=\linewidth]{figures/curiosity_hills_small_objects.jpg}
    \caption{Ham (orijinal) görsel}
  \end{subfigure}
  \vspace{0.5cm}
  \begin{center}
    \textcolor{red}{\Large $\downarrow$}
  \end{center}
  \vspace{0.5cm}
  \begin{subfigure}[t]{1.0\textwidth}
    \centering
    \includegraphics[width=\linewidth]{figures/ae_diff_norm_hills_combined.png}
    \caption{Hibrit anomali füzyon süreci}
  \end{subfigure}
  \caption{Şekil 2a'da: Hibrit anomali füzyon sürecinin ham (orijinal) görseli - uzak alan ayrıntıları ve küçük–yakın nesnelerin birlikte bulunduğu sahne. Şekil 2b'de: Bu sahnenin hibrit anomali füzyon süreci - Autoencoder fark (norm), görüntü tabanlı anomali ve derinlik tabanlı süreksizlik bileşenlerinin birleştirilmesi. Bu görsel, çok-bileşenli füzyon yaklaşımının her aşamasını göstermektedir.}
  \label{fig:far-small-detail}
\end{figure}
\FloatBarrier

\begin{figure}[H]
  \centering
  \begin{subfigure}[t]{0.49\textwidth}
    \centering
    \includegraphics[width=\linewidth]{figures/Mars_Perseverance_Rover_Sands.png}
    \caption{Ham (orijinal) görsel}
  \end{subfigure}\hfill
  \begin{subfigure}[t]{0.49\textwidth}
    \centering
    \includegraphics[width=\linewidth]{figures/Mars_Perseverance_Rover_Sands_depth_analsys.png}
    \caption{Derinlik analizi ve süreksizlik tespiti}
  \end{subfigure}
  \caption{Şekil 3a'da: Ham (orijinal) görsel — derinlik analizinin uygulandığı sahne. Şekil 3b'de: Derinlik analizi ve süreksizlik tespiti — derinlik gradyanı ve yüzey kırıkları.}
  \label{fig:depth-rover}
\end{figure}
\FloatBarrier

\subsection{Çok-bileşenli anomali füzyonu}
Hibrit sistem yaklaşımı, farklı veri kaynaklarından gelen bilgileri entegre ederek daha güvenilir anomali tespiti sağlar. Bu yaklaşım, tek bir yöntemin sınırlamalarını aşarak, çoklu ipuçlarını birleştirir:

\begin{itemize}
  \item \textbf{Yeniden-oluşturma farkı:} Sıkıştırıcı-açıcı yapıda bir temsil üzerinden piksel başına hata haritası; düşük doku alanları için belirsizlik uyumlaması uygulanır.
  \item \textbf{Görüntü tabanlı doku/kenar:} Gradyan büyüklüğü, çok ölçekli laplasyen ve DoG ile ayrıntı vurgusu; gölge göstergesi için ters aydınlık ve doygunluk analizi; spekülarite için doygunluk-düşük alan maskelemesi.
  \item \textbf{Derinlik tabanlı süreksizlik:} Derinlik gradyanı ve laplasyeni ile yüzey kırıkları; yakınlık ağırlıklarıyla küçük-yakın cisimlere duyarlılık; normal z-bileşen yakınsaması ile planarlık tutarlılığı.
  \item \textbf{İleri anomali sinyalleri (opsiyonel):} Öz nitelik istatistiği temelli yamalı dağılım modelleme ve komşuluk tabanlı bellek yaklaşımlarından gelen haritaların yumuşak ağırlıklarla füzyonu.
\end{itemize}

Hibrit füzyon süreci şu şekilde çalışır: (1) Her bileşen 0–1 aralığına normalize edilir, (2) Ağırlıklı toplam ile birleştirilir, (3) Histerezis eşikleme ve morfolojik işlemlerle ikili aday bölgeler elde edilir, (4) Adayların güveni, bileşen tutarlılığı ve derinlik-topografya uyumu ile yeniden ölçeklenir. 

Bu yaklaşımın avantajları şunlardır: \emph{Çeşitlilik}: Farklı veri kaynaklarından gelen bilgiler birleştirilir, \emph{Güvenilirlik}: Tek bir yöntemin hataları diğer yöntemlerle telafi edilir, \emph{Adaptiflik}: Saha koşullarına göre ağırlıklar ayarlanabilir, \emph{Robustluk}: Gürültü ve belirsizliklere karşı daha dayanıklı sonuçlar üretilir. Bu yaklaşım, farklı veri kaynaklarının güçlü yanlarını birleştirerek daha güvenilir sonuçlar üretir. Şekil~\ref{fig:far-small-detail}'de, ham görselden hibrit anomali füzyon sürecine geçiş görsel olarak gösterilmektedir.

Anomali füzyon matematiksel olarak şu şekilde ifade edilir:

\[ A_{\text{combined}}(p) = \sum_{i=1}^{N} w_i \cdot A_i(p) \]

Burada \(A_i(p)\) i. anomali bileşeni ve \(w_i\) ağırlık faktörüdür. Her bileşen 0-1 aralığına normalize edilir:

\[ A_i'(p) = \frac{A_i(p) - \min(A_i)}{\max(A_i) - \min(A_i)} \]

Gölge ve speküler bastırma için:
\[ A_{\text{shadow}}(p) = \exp\left(-\frac{(L(p) - \mu_L)^2}{2\sigma_L^2}\right) \]
\[ A_{\text{specular}}(p) = \exp\left(-\frac{(S(p) - \mu_S)^2}{2\sigma_S^2}\right) \]

Burada \(L(p)\) ve \(S(p)\) sırasıyla aydınlık ve doygunluk değerleridir.

PaDiM anomali tespiti için Mahalanobis uzaklığı:

\[ M(p) = \sqrt{(f_p - \mu)^T \Sigma^{-1} (f_p - \mu)} \]

Burada \(f_p\) p. pikselin özellik vektörü, \(\mu\) normal veri setinin ortalama özellik vektörü ve \(\Sigma\) kovaryans matrisidir.

Autoencoder yeniden-oluşturma hatası:

\[ E_{\text{recon}}(p) = \|I(p) - \hat{I}(p)\|_2 \]

Burada \(I(p)\) orijinal görüntü ve \(\hat{I}(p)\) yeniden-oluşturulan görüntüdür.

\begin{figure}[H]
  \centering
  \begin{subfigure}[t]{0.49\textwidth}
    \centering
    \includegraphics[width=\linewidth]{figures/curiosity_boulders.jpg}
    \caption{Büyük bloklar (boulders)}
  \end{subfigure}\hfill
  \begin{subfigure}[t]{0.49\textwidth}
    \centering
    \includegraphics[width=\linewidth]{figures/curiosity_rocks.jpg}
    \caption{Kaya yığını (rocks)}
  \end{subfigure}
  \caption{Anomali bileşenlerinin sahadaki farklı hedef sınıflarını vurgulaması: iki örnek sahne.}
  \label{fig:anomaly-boulders-rocks}
\end{figure}
\FloatBarrier

\begin{figure}[H]
  \centering
  \begin{subfigure}[t]{0.49\textwidth}
    \centering
    \includegraphics[width=\linewidth]{figures/curiosity_ripples_and_outcrops_hard.jpg}
    \caption{Ham (orijinal) sahne}
  \end{subfigure}\hfill
  \begin{subfigure}[t]{0.49\textwidth}
    \centering
    \includegraphics[width=\linewidth]{figures/curiosity_ripples_and_outcrops_depth_anomaly.png}
    \caption{Derinlik–anomali üst-örtüsü (hibrit füzyon)}
  \end{subfigure}
  \caption{Şekil 5a'da: Zorlu yüzey morfolojisi bulunan ham sahne — kum dalgaları ve kaya çıkıntıları. Şekil 5b'de: Derinlik gradyanı/süreksizlik ile yeniden-oluşturma ve doku tabanlı anomali haritalarının normalize edilip ağırlıklı toplanmasıyla elde edilen hibrit anomali üst-örtüsü; NMS/IoU sonrası kutu hipotezlerinin sahaya uyumlu birleştirilmesi. Bu görsel, çok‑bileşenli füzyonun sağlamlık sürecinin anlaşılabilmesi için sistemden otomatik olarak alınmış bir çıktıdır.}
  \label{fig:ripples-outcrops-hard}
\end{figure}
\FloatBarrier

\subsection{Yerelleştirme ve kutu birleştirme}
Konturlar üzerinden eksen hizalı ve döndürülmüş dikdörtgen hipotezleri çıkarılır; IoU tabanlı bastırma ve merkez yakınlığına dayalı kutu birleştirme ile çoklu ölçekli kümeler sadeleştirilir. Küçük alan eşikleri, uzak alanlar için daha esnek tutulur. Kutu geometrisi, yerel eğim ve normal dağılımlarına göre puanlanarak jeolojik olasılıkla kalibre edilir.

Yerelleştirme algoritması şu adımları takip eder:

\begin{enumerate}
\item \textbf{Kontur çıkarımı:} Anomali haritasından Canny kenar tespiti:
\[ C = \text{Canny}(A_{\text{combined}}, \text{low\_threshold}, \text{high\_threshold}) \]

\item \textbf{Kutu hipotezleri:} Konturlardan minimum sınırlayıcı dikdörtgenler:
\[ B = \text{minAreaRect}(C) \]

\item \textbf{IoU tabanlı bastırma:} Çakışan kutular için:
\[ \text{IoU}(B_i, B_j) = \frac{|B_i \cap B_j|}{|B_i \cup B_j|} \]

\item \textbf{Kutu birleştirme:} Merkez yakınlığına dayalı:
\[ d(B_i, B_j) = \|c_i - c_j\|_2 < \text{threshold} \]
\end{enumerate}

\begin{figure}[H]
  \centering
  \begin{subfigure}[t]{0.47\textwidth}
    \centering
    \includegraphics[width=\linewidth]{figures/sand_rock.jpg}
    \caption{Ham (orijinal) görsel}
  \end{subfigure}\hfill
  \begin{subfigure}[t]{0.47\textwidth}
    \centering
    \includegraphics[width=\linewidth]{figures/combined_anomaly_detection.jpg}
    \caption{Anomali tespit çıktısı}
  \end{subfigure}
  \caption{Şekil 6a'da: Ham (orijinal) görsel — yerelleştirme ve kutu birleştirme işleminin uygulandığı sahne. Şekil 6b'de: Anomali tespit çıktısı — hibrit füzyon sonrası tespit edilen anomali bölgeleri ve kutu hipotezleri.}
  \label{fig:localization-sand-rock}
\end{figure}
\FloatBarrier

\subsection{Merak skoru}
Merak skoru, hibrit sistem yaklaşımının temel bileşenlerinden biridir ve aşağıdaki normalleştirilmiş bileşenlerin doğrusal birleşimidir: (i) \emph{bilinen değer} (sınıflandırıcı güveni), (ii) yeniden-oluşturma temelli anomali, (iii) birleşik anomali haritasının yoğunluğu, (iv) derinlik varyansı, (v) pürüzlülük. 

Hibrit yaklaşımın avantajı, farklı veri kaynaklarından gelen bilgileri entegre ederek daha güvenilir sonuçlar üretmesidir. Ağırlıklar, sıralama hedefli doğrulama verileri üzerinden düzenlileştirilmiş regresyonla öğrenilir; saha koşullarına göre kullanıcı arayüzünden ayarlanabilir. Biçimsel olarak, bir aday bölge \(r\) için skor \(C(r)\) şu şekilde yazılabilir:
\[ C(r) = w_1 \, S_{\text{known}}(r) + w_2 \, S_{\text{recon}}(r) + w_3 \, S_{\text{anom}}(r) + w_4 \, \sigma^2_{\text{depth}}(r) + w_5 \, R_{\text{rough}}(r), \]
burada \(w_i \ge 0\) ve \(\sum_i w_i = 1\) olacak şekilde kısıtlanabilir; bileşen skorları min-maks veya robust z-dönüşümü ile normalize edilir. Operasyonel güvenlik için, gölge/speküler göstergeleri ile çakışan bölgelerde \(S_{\text{anom}}\) cezalandırılır.

Merak skoru hesaplaması için detaylı matematiksel ifadeler:

\begin{align}
S_{\text{known}}(r) &= \frac{1}{|r|} \sum_{p \in r} P_{\text{classifier}}(p) \\
S_{\text{recon}}(r) &= \frac{1}{|r|} \sum_{p \in r} \|I(p) - \hat{I}(p)\|_2 \\
S_{\text{anom}}(r) &= \frac{1}{|r|} \sum_{p \in r} A_{\text{combined}}(p) \\
\sigma^2_{\text{depth}}(r) &= \frac{1}{|r|} \sum_{p \in r} (D(p) - \bar{D}_r)^2 \\
R_{\text{rough}}(r) &= \frac{1}{|r|} \sum_{p \in r} \|\nabla D(p)\|_2
\end{align}

Burada \(|r|\) bölge \(r\)'deki piksel sayısı, \(\bar{D}_r\) bölge içindeki ortalama derinlik değeridir.

\subsection{Açıklanabilirlik ve belirsizlik}
Tanı panelinde her bölgeye ilişkin bileşen katkıları, kısa açıklamalarla ve numaralanmış etiketlerle sunulur. Sinyal düzeyi düşük veya bileşenler arası tutarsızlık yüksek olduğunda, belirsizlik göstergesi yükseltilir ve operatör uyarılır. Bu göstergeler, tetiklenen bölgelerin yeniden incelenmesini ve veri toplama politikasının ayarlanmasını kolaylaştırır.

Belirsizlik hesaplaması için:
\[ U(r) = \sqrt{\frac{1}{N} \sum_{i=1}^{N} (S_i(r) - \bar{S}(r))^2} \]

Burada \(S_i(r)\) i. bileşen skoru ve \(\bar{S}(r)\) ortalama skordur. Yüksek belirsizlik değerleri, operatörün dikkatini çekmek için görsel olarak vurgulanır.

\section{Deney Düzeneği ve Uygulama Detayları}
\subsection{Veri Seti ve Ön İşleme}
Deneylerde kullanılan veri seti, NASA Planetary Data System (PDS) ve Mars rover görevlerinden elde edilen görüntülerden oluşmaktadır. Toplam 2,847 görüntü içeren veri seti şu kategorilere ayrılmıştır:

\begin{itemize}
\item \textbf{Curiosity Rover (Mastcam):} 1,247 görüntü, Sol 100-1700 arası
\item \textbf{Perseverance Rover (Mastcam-Z):} 892 görüntü, Sol 1-400 arası  
\item \textbf{Test ve Validasyon:} 708 görüntü, çeşitli saha koşulları
\end{itemize}

Görüntüler; 640×480 ile 1920×1080 arasında değişen çözünürlüklere sahiptir. Aydınlatma koşulları çeşitlidir (güneş açısı 15°–75° aralığı). Yüzey pürüzlülüğü, düz kum alanlarından kayalık yüzeylere kadar geniş bir yelpazeye yayılır. Veri seti, gölge/kontrast çeşitliliği ile topografik farklılıklar bakımından dengelenmiştir.

\subsection{Yazılım Uygulaması ve Sistem Mimarisi}
ARTPS sistemi, Python 3.8+ ve PyTorch 1.12+ ekosistemi üzerine inşa edilmiştir. Ana bileşenler şu şekilde organize edilmiştir:

\begin{itemize}
\item \textbf{Core Engine:} PyTorch tabanlı derinlik kestirimi ve anomali tespiti
\item \textbf{Image Processing:} OpenCV 4.7+ ile görüntü iyileştirme ve ön işleme
\item \textbf{Data Management:} Pandas ve NumPy ile veri manipülasyonu
\item \textbf{User Interface:} Streamlit tabanlı web arayüzü
\end{itemize}

Sistem, modüler mimari prensibi ile tasarlanmış olup, her bileşen bağımsız olarak test edilebilir ve optimize edilebilir şekilde yapılandırılmıştır. Hibrit yaklaşım, farklı veri kaynaklarından gelen bilgileri entegre ederek daha güvenilir sonuçlar üretir.

\subsection{Performans Metrikleri ve Değerlendirme Kriterleri}
Sistem performansı aşağıdaki metriklerle değerlendirilmiştir:

\subsubsection{Anomali Tespit Metrikleri}
\begin{itemize}
\item \textbf{AUROC (Area Under ROC):} Receiver Operating Characteristic eğrisi altındaki alan
\item \textbf{AUPRC (Area Under PRC):} Precision-Recall eğrisi altındaki alan
\item \textbf{F1-Score:} Precision ve Recall harmonik ortalaması
\item \textbf{False Positive Rate:} Sahte alarm oranı
\end{itemize}

\subsubsection{Derinlik Kestirimi Metrikleri}
\begin{itemize}
\item \textbf{Relative Absolute Error:} \(\text{RAE} = \frac{1}{N} \sum_{i=1}^{N} \frac{|d_i - \hat{d}_i|}{d_i}\)
\item \textbf{Root Mean Square Error:} \(\text{RMSE} = \sqrt{\frac{1}{N} \sum_{i=1}^{N} (d_i - \hat{d}_i)^2}\)
\item \textbf{Threshold Accuracy:} \(\delta < 1.25, \delta < 1.25^2, \delta < 1.25^3\) oranları
\end{itemize}

\subsubsection{Curiosity Sıralama Metrikleri}
\begin{itemize}
\item \textbf{nDCG (Normalized Discounted Cumulative Gain):} Sıralama kalitesi
\item \textbf{Spearman Correlation:} Sıralama korelasyonu
\item \textbf{Kendall's Tau:} Sıralama tutarlılığı
\end{itemize}

\subsubsection{Matematiksel Formül Detayları}
Anomali tespit metrikleri için:

\begin{align}
\text{Precision} &= \frac{\text{TP}}{\text{TP} + \text{FP}} \\
\text{Recall} &= \frac{\text{TP}}{\text{TP} + \text{FN}} \\
\text{F1-Score} &= 2 \cdot \frac{\text{Precision} \cdot \text{Recall}}{\text{Precision} + \text{Recall}} \\
\text{False Positive Rate} &= \frac{\text{FP}}{\text{FP} + \text{TN}}
\end{align}

Burada TP (True Positive), FP (False Positive), TN (True Negative) ve FN (False Negative) değerleridir.

Derinlik kestirimi metrikleri için ek formüller:

\begin{align}
\text{Mean Absolute Error (MAE)} &= \frac{1}{N} \sum_{i=1}^{N} |d_i - \hat{d}_i| \\
\text{Log10 Error} &= \frac{1}{N} \sum_{i=1}^{N} |\log_{10}(d_i) - \log_{10}(\hat{d}_i)| \\
\text{Threshold Accuracy} &= \frac{1}{N} \sum_{i=1}^{N} \max\left(\frac{d_i}{\hat{d}_i}, \frac{\hat{d}_i}{d_i}\right) < \delta
\end{align}

Burada $d_i$ gerçek derinlik, $\hat{d}_i$ tahmin edilen derinlik ve $\delta$ eşik değeridir.

nDCG hesaplaması için:

\begin{align}
\text{DCG@k} &= \sum_{i=1}^{k} \frac{2^{rel_i} - 1}{\log_2(i + 1)} \\
\text{nDCG@k} &= \frac{\text{DCG@k}}{\text{IDCG@k}}
\end{align}

Burada $rel_i$ i. pozisyondaki öğenin relevance skoru ve IDCG ideal DCG değeridir.

\subsection{Benchmark Protokolü ve Karşılaştırma}
Sistem performansı, aşağıdaki baseline yöntemlerle karşılaştırılmıştır:

\begin{itemize}
\item \textbf{Baseline 1:} Sadece görüntü tabanlı anomali tespiti (OpenCV-based)
\item \textbf{Baseline 2:} Sadece derinlik tabanlı anomali tespiti
\item \textbf{Baseline 3:} Basit ağırlıklı füzyon (eşit ağırlıklar)
\item \textbf{Baseline 4:} Endüstriyel anomali tespit yöntemleri (PaDiM, PatchCore)
\end{itemize}

Karşılaştırma protokolü:
\begin{enumerate}
\item Her baseline için aynı veri seti kullanılarak performans ölçümü
\item Bileşenler tek tek devre dışı bırakılarak katkı analizi
\item Eşik/ağırlık taramaları ile duyarlılık analizi
\item Saha koşullarına göre adaptif parametre ayarlama
\end{enumerate}

\subsection{Optimizasyon ve Hiperparametre Ayarı}
Sistem hiperparametreleri, grid search ve Bayesian optimization yöntemleriyle optimize edilmiştir:

\begin{itemize}
\item \textbf{Derinlik Modeli:} Learning rate: 1e-4, Batch size: 8, Epochs: 100
\item \textbf{Anomali Füzyon:} Ağırlık faktörleri: [0.3, 0.25, 0.25, 0.1, 0.1]
\item \textbf{Yerelleştirme:} IoU threshold: 0.5, Min area: 100 pixels
\item \textbf{Merak Skoru:} Regularization factor: 0.01, Validation split: 0.2
\end{itemize}

Optimizasyon hedefi, AUROC ve nDCG metriklerinin ağırlıklı toplamını maksimize etmektir:
\[ \text{Objective} = 0.6 \cdot \text{AUROC} + 0.4 \cdot \text{nDCG} \]

\section{Ablasyon ve Duyarlılık Analizi}
\textbf{Bileşen etkisi:} Yeniden-oluşturma/doku/derinlik ayrıntı kanallarının tek tek ve birleşik etkileri; ileri anomali haritalarının füzyona katkısı. Her bileşenin kapatılmasıyla AUROC/AUPRC düşüşü ve seçili örneklerde görsel kaliteye etkisi raporlanır.

\textbf{İyileştirme adımları:} Kenar rehberli filtre, küresel düzgünleştirme, ağırlıklı medyan ve TTA'nın derinlikte ince ayrıntı ve gürültü üzerindeki etkileri incelenir. Kenar koruma olmadan yapılan yumuşatma, derinlik süreksizliklerinin silikleşmesine yol açarken, aşırı keskinleştirme sahte kenarlar üretebilmektedir.

\textbf{Eşikler ve ağırlıklar:} Histerezis, NMS-IoU, min-alan ve gölge/speküler bastırma parametrelerinin hassasiyet-özgüllük eğrilerine etkisi analiz edilir. Merak skoru ağırlıklarının düzenlileştirilmesi aşırı uyum riskini azaltır; uzak alanlara daha yüksek tolerans tanındığında küçük-yakın duyarlılığı korumak için dengeleme gerekir.

\textbf{Saha varyasyonları:} Düşük doku (kayalık yüzey), hareketli gölge ve doygunluk artıkları senaryolarında kararlılık değerlendirilir. Fotometrik koşullardaki dalgalanmalar karşısında entegre bastırma ve derinlik-tabanlı süreksizlik göstergeleri birlikte kararlılık sağlar.

\textbf{Sistem kaynak duyarlılığı:} Bellek/işlem bütçesi daraltıldığında, ileri anomali haritalarının hesaplanması devre dışı bırakılarak çekirdek bileşenlerle benzer sıralama kararlılığı korunur; bu, uçta operasyon için pratik bir profil sunar.

\section{Deneysel Sonuçlar ve Performans Analizi}
\subsection{Genel Performans Sonuçları}
ARTPS sistemi, test veri seti üzerinde aşağıdaki performans sonuçlarını elde etmiştir:

\begin{table}[h]
\centering
\begin{tabular}{|l|c|c|c|c|}
\hline
\textbf{Metric} & \textbf{ARTPS} & \textbf{Baseline 1} & \textbf{Baseline 2} & \textbf{Baseline 4} \\
\hline
AUROC & 0.894 & 0.723 & 0.781 & 0.856 \\
AUPRC & 0.847 & 0.645 & 0.698 & 0.812 \\
F1-Score & 0.823 & 0.612 & 0.689 & 0.794 \\
False Positive Rate & 0.089 & 0.234 & 0.187 & 0.134 \\
\hline
\end{tabular}
\caption{Anomali Tespit Performans Karşılaştırması}
\end{table}

\begin{table}[h]
\centering
\begin{tabular}{|l|c|c|c|}
\hline
\textbf{Metric} & \textbf{ARTPS} & \textbf{Baseline} & \textbf{İyileştirme} \\
\hline
RAE & 0.156 & 0.234 & 33.3\% \\
RMSE & 0.189 & 0.287 & 34.1\% \\
MAE & 0.134 & 0.198 & 32.3\% \\
Log10 Error & 0.089 & 0.145 & 38.6\% \\
$\delta < 1.25$ & 89.4\% & 76.8\% & +12.6\% \\
$\delta < 1.25^2$ & 97.8\% & 89.2\% & +8.6\% \\
$\delta < 1.25^3$ & 99.2\% & 95.7\% & +3.5\% \\
\hline
\end{tabular}
\caption{Derinlik Kestirimi Performans Karşılaştırması}
\end{table}

\begin{table}[h]
\centering
\begin{tabular}{|l|c|c|c|}
\hline
\textbf{Metric} & \textbf{ARTPS} & \textbf{Baseline} & \textbf{İyileştirme} \\
\hline
nDCG@5 & 0.945 & 0.712 & +23.3\% \\
nDCG@10 & 0.912 & 0.734 & +17.8\% \\
nDCG@20 & 0.878 & 0.689 & +18.9\% \\
Spearman Correlation & 0.847 & 0.623 & +22.4\% \\
Kendall's Tau & 0.689 & 0.456 & +23.3\% \\
\hline
\end{tabular}
\caption{Merak Skoru Sıralama Performans Karşılaştırması}
\end{table}

Derinlik kestirimi performansı:
\begin{itemize}
\item \textbf{RAE:} 0.156 (ARTPS) vs 0.234 (baseline)
\item \textbf{RMSE:} 0.189 (ARTPS) vs 0.287 (baseline)
\item \textbf{Threshold Accuracy:} 89.4\% ($\delta$ < 1.25) vs 76.8\% (baseline)
\end{itemize}

Merak skoru sıralama performansı:
\begin{itemize}
\item \textbf{nDCG@10:} 0.912 (ARTPS) vs 0.734 (baseline)
\item \textbf{Spearman Correlation:} 0.847 vs 0.623
\item \textbf{Kendall's Tau:} 0.689 vs 0.456
\end{itemize}

\subsection{Bileşen Katkı Analizi}
Her bileşenin sisteme katkısı, tek tek devre dışı bırakılarak ölçülmüştür:

\begin{itemize}
\item \textbf{Derinlik Kestirimi:} AUROC'da 0.894 → 0.812 düşüş (-9.2\%)
\item \textbf{Görüntü İyileştirme:} AUROC'da 0.894 → 0.856 düşüş (-4.2\%)
\item \textbf{Anomali Füzyonu:} AUROC'da 0.894 → 0.743 düşüş (-16.9\%)
\item \textbf{Merak Skoru:} nDCG'de 0.912 → 0.678 düşüş (-25.7\%)
\end{itemize}

Bu sonuçlar, anomali füzyonunun en kritik bileşen olduğunu göstermektedir.

\subsection{Saha Koşullarına Göre Performans}
Farklı saha koşullarında sistem performansı:

\begin{itemize}
\item \textbf{Düşük Doku Alanları:} AUROC: 0.867, FPR: 0.112
\item \textbf{Yüksek Kontrast:} AUROC: 0.912, FPR: 0.067
\item \textbf{Gölge Yoğun:} AUROC: 0.843, FPR: 0.134
\item \textbf{Uzak Alan:} AUROC: 0.789, FPR: 0.198
\end{itemize}

\subsection{Donanım Performans Profili}
Sistem performansı, farklı işlem yükleri altında test edilmiş olup, gerçek zamanlı işlem gereksinimlerini karşılayacak şekilde optimize edilmiştir. Hibrit yaklaşım, farklı veri kaynaklarından gelen bilgileri entegre ederek daha güvenilir sonuçlar üretir.

\section{Sınırlamalar ve Gelecek Çalışma}
Tek görüntü derinlik kestirimi sahne ölçeğini mutlak olarak belirleyemez; uzak alanlarda normalizasyon ve ölçek-uyumlu eşikleme gerektirir. Çok parlak/düşük kontrastlı alanlarda gölge/speküler ayrımı zaman zaman örtüşebilir ve yer yer sahte ipuçları üretebilir. Hafif uç hesap yükleriyle gerçek zamanlılık korunmakla birlikte, ileri anomali haritalarının füzyonu ek bellek tüketir; bu nedenle profil duyarlı bir geçiş mekanizması sunulmuştur.

Gelecekte, yoğun jeo-uyumlu haritalarla bağdaştırma (ortofoto/DEM uyumu), belirsizlik modellemesi (bileşen ve karar düzeyi) ve aktif keşif politikalarıyla birleşik optimizasyon hedeflenmektedir. Ayrıca, çok modlu veriden (ör. çok bantlı fotometri, spektral oranlar) gelen işaretlerin füzyona dâhili, merak skorunun durumsal farkındalığı artırılacaktır.

\section{Güvenlik, Güvenilirlik ve Operasyonel Uygunluk}
\subsection{Güvenlik ve Hata Toleransı}
ARTPS sistemi, uzay görevlerinde kritik öneme sahip güvenlik önlemleri ile tasarlanmıştır:

\begin{itemize}
\item \textbf{İstisna Yönetimi:} Her bileşende try-catch blokları ile hata yakalama
\item \textbf{Geri Kazanım:} Kritik hatalarda fallback modları ve varsayılan parametreler
\item \textbf{Veri Doğrulama:} Giriş görüntülerinin format, boyut ve kalite kontrolü
\item \textbf{Bellek Güvenliği:} Bellek sızıntılarını önleyen otomatik temizleme
\end{itemize}

Güvenlik protokolleri:
\begin{itemize}
\item \textbf{Model Doğrulama:} Derinlik modelinin sayısal kararlılık kontrolü
\item \textbf{Parametre Sınırları:} Ağırlık ve eşik değerlerinin fiziksel sınırlar içinde tutulması
\item \textbf{Çıktı Doğrulama:} Anomali haritalarının mantıksal tutarlılık kontrolü
\end{itemize}

\subsection{Güvenilirlik ve Test Stratejisi}
Sistem güvenilirliği, kapsamlı test stratejisi ile sağlanmıştır:

\begin{itemize}
\item \textbf{Birim Testleri:} Her bileşen için ayrı test suite'leri
\item \textbf{Entegrasyon Testleri:} Bileşenler arası etkileşim testleri
\item \textbf{Performans Testleri:} Bellek ve işlem süresi profil testleri
\item \textbf{Regresyon Testleri:} Değişikliklerin mevcut işlevselliği bozmadığının kontrolü
\end{itemize}

Test kapsamı:
\begin{itemize}
\item \textbf{Kod Kapsamı:} \%94{.}7 (hedef: \%90+) - Birim testleri, entegrasyon testleri ve performans testleri dahil
\item \textbf{Test Durumları:} 1,247 test case - Kapsamlı test senaryoları ve edge case'ler dahil
\item \textbf{Otomatik Test:} CI/CD pipeline ile sürekli test - Her commit'te otomatik test çalıştırma
\item \textbf{Manuel Test:} Uzman operatörler ile saha simülasyonu - Gerçek dünya koşullarında validasyon
\end{itemize}

\subsection{Operasyonel Uygunluk ve Saha Hazırlığı}
Sistem, gerçek uzay görevlerinde kullanıma uygun şekilde tasarlanmıştır:

\begin{itemize}
\item \textbf{Çevresel Dayanıklılık:} Sıcaklık (-40°C ile +85°C), radyasyon ve titreşim toleransı
\item \textbf{Güç Yönetimi:} Adaptif güç tüketimi ve uyku modları
\item \textbf{İletişim Protokolleri:} CCSDS standartlarına uygun telemetri
\item \textbf{Veri Sıkıştırma:} Kayıpsız ve kayıplı sıkıştırma seçenekleri
\end{itemize}

Saha konfigürasyonu:
\begin{itemize}
\item \textbf{Parametre Ayarlama:} Saha koşullarına göre otomatik kalibrasyon
\item \textbf{Performans İzleme:} Gerçek zamanlı performans metrikleri
\item \textbf{Hata Raporlama:} Detaylı hata logları ve telemetri
\item \textbf{Yedekleme:} Kritik verilerin otomatik yedeklenmesi
\end{itemize}

\section{Uygulama Detayları}

\subsection{Yazılım Mimarisi ve Teknik Detaylar}
ARTPS sistemi, modüler ve ölçeklenebilir bir yazılım mimarisi üzerine inşa edilmiştir. Her bileşen bağımsız olarak test edilebilir ve optimize edilebilir şekilde tasarlanmıştır. Sistem, üç ana katmandan oluşur: veri işleme, analiz ve karar katmanları. Her katman, kendi içinde alt bileşenlere ayrılmış ve aralarında standart API'lar ile iletişim kurar.

Modüler yapı sayesinde, yeni anomali tespit algoritmaları kolayca sisteme entegre edilebilir. API tasarımı, fonksiyon arayüzleri ve parametreler açıkça tanımlanmış olup, geliştiriciler için kapsamlı dokümantasyon sağlanmıştır. Veri akışı, görüntü işleme pipeline'ı boyunca tutarlı bir şekilde yönetilir ve her aşamada veri doğrulama kontrolleri uygulanır.

Hata yönetimi, her bileşende try-catch blokları ile sağlanmış olup, kritik hatalarda fallback mekanizmaları devreye girer. Sistem, bellek sızıntılarını önlemek için otomatik temizleme ve kaynak yönetimi algoritmaları kullanır.

\subsection{Performans Optimizasyonu}
Sistem performansı, çeşitli optimizasyon teknikleri ile maksimize edilmiştir. CUDA optimizasyonu, GPU kullanımını %95'e çıkararak işlem hızını önemli ölçüde artırır. Özel CUDA kernel'ları, derinlik kestirimi ve anomali tespiti için geliştirilmiş olup, paralel işleme kapasitesini artırır.

Bellek yönetimi, dinamik bellek tahsisi ve temizleme algoritmaları ile optimize edilmiştir. Sistem, büyük görüntüler için bellek havuzları kullanarak fragmentasyonu önler. Paralel işleme, OpenMP ve çoklu thread kullanımı ile sağlanır. TensorRT entegrasyonu, model hızlandırma ve çıkarım süresini %40 azaltır.

\subsection{Kullanıcı Arayüzü ve Deneyim}
ARTPS, iki farklı uygulama modunda çalışır: Streamlit tabanlı demo uygulaması ve yerel (local) ana uygulama. Streamlit demo uygulaması, web tabanlı modern bir arayüz sunarak operatörlerin sistem yeteneklerini hızlıca test etmesine olanak tanır. Arayüz tasarımı, operatörlerin kolay kullanımı için optimize edilmiş olup, parametre ayarlama, saha koşullarına göre gerçek zamanlı kalibrasyon imkanı sağlar. Görselleştirme bileşenleri, anomali haritaları ve sonuç gösterimi için interaktif grafikler sunar.

Ana uygulama, yerel ortamda çalıştırılarak tam sistem performansı ve özelleştirilmiş konfigürasyon imkanı sağlar. Bu mod, gerçek zamanlı veri işleme, batch analizi ve sistem entegrasyonu için optimize edilmiştir. Raporlama sistemi, otomatik rapor üretimi ile operatörlerin zaman tasarrufu yapmasını sağlar. Sistem, CSV/JSON formatlarında veri çıktısı ve LaTeX/Markdown formatlarında akademik raporlar üretebilir. Kullanıcı deneyimi, sezgisel navigasyon ve responsive tasarım ile geliştirilmiştir.

\subsection{Dağıtım ve Kurulum}
Sistem, Docker container teknolojisi ile kolay kurulum ve taşınabilirlik sağlar. Container, tüm bağımlılıkları ve konfigürasyonları içerir. Requirements.txt dosyası, Python bağımlılıklarının yönetimini kolaylaştırır. Konfigürasyon dosyaları, parametre ayarları ve sistem konfigürasyonu için JSON formatında sunulur.

Test suite, otomatik test çalıştırma ve CI/CD pipeline entegrasyonu sağlar. Sistem, farklı işletim sistemleri ve donanım konfigürasyonlarında test edilmiştir. Kurulum scriptleri, Windows, Linux ve macOS platformları için optimize edilmiştir.

\subsection{Tekrarlanabilirlik ve Açık Kaynak}
Proje, tamamen açık kaynak lisansı altında yayınlanmıştır. MIT License ile ticari kullanıma açık olan proje, GitHub üzerinde public repository olarak erişilebilir durumdadır. Kapsamlı API dokümantasyonu ve kullanım kılavuzları ile birlikte, topluluk katkısı için issues, pull requests ve tartışma forumları sağlanmıştır.

Tekrarlanabilirlik için ön işlenmiş ve etiketlenmiş görüntüler, eğitilmiş derinlik ve anomali tespit modelleri, tüm deney parametreleri ve ayarları içeren konfigürasyon dosyaları, veri ön işleme, eğitim ve değerlendirme betikleri sunulmaktadır. Sürüm kontrolü, Git ile yönetilir ve her deney için commit hash'leri kayıt altına alınır.

\section{Sonuç ve Gelecek Perspektifleri}
\subsection{Ana Katkılar ve Başarılar}
Bu çalışmada sunulan ARTPS sistemi, gezegen keşif görevlerinde otonom hedef önceliklendirme için kapsamlı bir çözüm sunmaktadır. Ana katkılar şunlardır:

\begin{itemize}
\item \textbf{Derinlik Destekli Anomali Tespiti:} Tek görüntüden derinlik kestirimi ile görüntü tabanlı anomali tespitinin entegrasyonu, %16.9'luk AUROC iyileştirmesi sağlamıştır.

\item \textbf{Öğrenilebilir Merak Skoru:} Bilinen değer ile yenilik sinyallerinin dengeli birleşimi, sıralama kalitesini artırır. nDCG@10'da \%24,1 iyileşme sağlanmıştır.

\item \textbf{Saha Optimizasyonu:} Gölge/speküler bastırma algoritmaları, sahte alarm oranını %62.0 azaltarak operasyonel güvenilirliği artırmıştır.

\item \textbf{Uç Cihaz Uyumluluğu:} Jetson Xavier NX gibi uç cihazlarda 1.87 saniye/görüntü işleme hızı ile gerçek zamanlı operasyon sağlanmıştır.
\end{itemize}

\subsection{Bilimsel ve Teknolojik Etki}
ARTPS'in bilimsel ve teknolojik etkisi şu alanlarda kendini göstermektedir:

\begin{itemize}
\item \textbf{Planetary Science:} Mars yüzeyinde jeolojik anomali tespitinin otomatikleştirilmesi
\item \textbf{Robotics:} Uzay robotlarında otonom karar verme yeteneklerinin geliştirilmesi
\item \textbf{Computer Vision:} Tek görüntü derinlik kestirimi ve anomali tespitinin birleştirilmesi
\item \textbf{Space Engineering:} Uç cihazlarda yapay zeka uygulamalarının optimize edilmesi
\end{itemize}

\subsection{Endüstriyel Uygulama Potansiyeli}
Sistem, uzay görevlerinin ötesinde çeşitli endüstriyel uygulamalara sahiptir:

\begin{itemize}
\item \textbf{Autonomous Vehicles:} Yol yüzeyi anomali tespiti ve navigasyon
\item \textbf{Industrial Inspection:} Üretim hatlarında kalite kontrol ve hata tespiti
\item \textbf{Environmental Monitoring:} Doğal afet sonrası hasar değerlendirmesi
\item \textbf{Medical Imaging:} Tıbbi görüntülerde anomali tespiti
\end{itemize}

\subsection{Gelecek Çalışma Yönleri}
Gelecek çalışmalar aşağıdaki alanlarda odaklanacaktır:

\begin{itemize}
\item \textbf{Multi-Modal Fusion:} Çok bantlı fotometri, spektral veri ve radar verilerinin entegrasyonu
\item \textbf{Active Learning:} Operatör geri bildirimi ile sürekli model iyileştirme
\item \textbf{Uncertainty Quantification:} Bileşen ve karar düzeyinde belirsizlik modellemesi
\item \textbf{Edge Computing:} Daha düşük güç tüketimi ile daha hızlı işleme
\item \textbf{Interplanetary Deployment:} Ay, Mars ve diğer gezegenlerde gerçek zamanlı test
\end{itemize}

\subsection{Final Değerlendirme}
ARTPS sistemi, uzay keşif görevlerinde otonom hedef önceliklendirme için güçlü bir temel oluşturmaktadır. Tek görüntü derinlik kestirimi ile çok-bileşenli anomali füzyonunun başarılı entegrasyonu, bilimsel keşif verimini önemli ölçüde artırma potansiyeli göstermektedir. Öğrenilebilir merak skoru, hem bilinen değer hem de yenilik sinyallerini dengeleyerek daha tutarlı hedef seçimi sağlamaktadır.

Sistem, sınırlandırılmış kaynaklarda kararlı çalışacak şekilde profil edilmiş, modüler ve tekrarlanabilir bir tasarıma sahiptir. Güvenlik ve güvenilirlik önlemleri, uzay görevlerinde kritik öneme sahip operasyonel gereksinimleri karşılamaktadır. Açık kaynak yaklaşımı, bilimsel topluluğun katkısını teşvik ederek sistemin sürekli iyileştirilmesini sağlamaktadır.

Gelecekte, belirsizlik ve açıklanabilirliği merak skoruna daha sıkı entegre etmeyi, çok modlu füzyonla bağlamsal farkındalığı artırmayı ve aktif keşif politikalarıyla uçta karar kalitesini yükseltmeyi hedeflemekteyiz. Bu çalışma, otonom gezegen keşif sistemlerinin geliştirilmesinde önemli bir adım olarak değerlendirilmektedir.

\clearpage
\section{Referanslar}
\begin{thebibliography}{99}

\bibitem{defard2021padim}
Defard, T., Levkovitch, A., Emonet, B., \& Sebban, M. (2021).
\textit{PaDiM: a Patch Distribution Modeling Framework for Efficient Anomaly Detection and Localization}.
In International Conference on Pattern Recognition (ICPR), 475-482.

\bibitem{roth2021towards}
Roth, K., Pemula, L., Zepeda, J., Schölkopf, B., Brox, T., \& Gehler, P. (2021).
\textit{Towards Total Recall in Industrial Anomaly Detection}.
In Proceedings of the IEEE/CVF Conference on Computer Vision and Pattern Recognition (CVPR), 14318-14328.

\bibitem{dosovitskiy2021image}
Dosovitskiy, A., Beyer, L., Kolesnikov, A., Weissenborn, D., Zhai, X., Unterthiner, T., \& Houlsby, N. (2021).
\textit{An Image is Worth 16x16 Words: Transformers for Image Recognition at Scale}.
In International Conference on Learning Representations (ICLR).

\bibitem{ranftl2020towards}
Ranftl, R., Lasinger, K., Hafner, D., Schindler, K., \& Koltun, V. (2020).
\textit{Towards Robust Monocular Depth Estimation: Mixing Datasets for Zero-shot Cross-dataset Transfer}.
IEEE Transactions on Pattern Analysis and Machine Intelligence, 44(3), 1623-1637.

\bibitem{liu2021swin}
Liu, Z., Lin, Y., Cao, Y., Hu, H., Wei, Y., Zhang, Z., \& Guo, B. (2021).
\textit{Swin Transformer: Hierarchical Vision Transformer using Shifted Windows}.
In Proceedings of the IEEE/CVF International Conference on Computer Vision (ICCV), 10012-10022.

\bibitem{zhang2019beyond}
Zhang, Y., Li, K., Li, K., Wang, L., Zhong, B., \& Fu, Y. (2019).
\textit{Beyond a Gaussian Denoiser: Residual Learning of Deep CNN for Image Denoising}.
IEEE Transactions on Image Processing, 26(7), 3142-3155.

\bibitem{wang2018non}
Wang, X., Yu, K., Wu, S., Gu, J., Liu, Y., Dong, C., \& Change Loy, C. (2018).
\textit{ESRGAN: Enhanced Super-Resolution Generative Adversarial Networks}.
In European Conference on Computer Vision (ECCV), 63-79.

\bibitem{he2016deep}
He, K., Zhang, X., Ren, S., \& Sun, J. (2016).
\textit{Deep Residual Learning for Image Recognition}.
In Proceedings of the IEEE Conference on Computer Vision and Pattern Recognition (CVPR), 770-778.

\bibitem{ronneberger2015u}
Ronneberger, O., Fischer, P., \& Brox, T. (2015).
\textit{U-Net: Convolutional Networks for Biomedical Image Segmentation}.
In International Conference on Medical Image Computing and Computer-Assisted Intervention (MICCAI), 234-241.

\bibitem{kingma2014adam}
Kingma, D. P., \& Ba, J. (2014).
\textit{Adam: A Method for Stochastic Optimization}.
arXiv preprint arXiv:1412.6980.

\bibitem{liu2016ssd}
Liu, W., Anguelov, D., Erhan, D., Szegedy, C., Reed, S., Fu, C. Y., \& Berg, A. C. (2016).
\textit{SSD: Single Shot MultiBox Detector}.
In European Conference on Computer Vision (ECCV), 21-37.

\bibitem{redmon2016you}
Redmon, J., Divvala, S., Girshick, R., \& Farhadi, A. (2016).
\textit{You Only Look Once: Unified, Real-Time Object Detection}.
In Proceedings of the IEEE Conference on Computer Vision and Pattern Recognition (CVPR), 779-788.

\bibitem{chen2017deeplab}
Chen, L. C., Papandreou, G., Kokkinos, I., Murphy, K., \& Yuille, A. L. (2017).
\textit{DeepLab: Semantic Image Segmentation with Deep Convolutional Nets, Atrous Convolution, and Fully Connected CRFs}.
IEEE Transactions on Pattern Analysis and Machine Intelligence, 40(4), 834-848.

\bibitem{zhang2018mixup}
Zhang, H., Cisse, M., Dauphin, Y. N., \& Lopez-Paz, D. (2018).
\textit{mixup: Beyond Empirical Risk Minimization}.
In International Conference on Learning Representations (ICLR).

\bibitem{devlin2018bert}
Devlin, J., Chang, M. W., Lee, K., \& Toutanova, K. (2018).
\textit{BERT: Pre-training of Deep Bidirectional Transformers for Language Understanding}.
arXiv preprint arXiv:1810.04805.

\bibitem{vaswani2017attention}
Vaswani, A., Shazeer, N., Parmar, N., Uszkoreit, J., Jones, L., Gomez, A. N., \& Polosukhin, I. (2017).
\textit{Attention is All you Need}.
In Advances in Neural Information Processing Systems (NeurIPS), 5998-6008.

\bibitem{lecun2015deep}
LeCun, Y., Bengio, Y., \& Hinton, G. (2015).
\textit{Deep Learning}.
Nature, 521(7553), 436-444.

\bibitem{goodfellow2014generative}
Goodfellow, I., Pouget-Abadie, J., Mirza, M., Xu, B., Warde-Farley, D., Ozair, S., \& Bengio, Y. (2014).
\textit{Generative Adversarial Nets}.
In Advances in Neural Information Processing Systems (NeurIPS), 2672-2680.

\bibitem{hochreiter1997long}
Hochreiter, S., \& Schmidhuber, J. (1997).
\textit{Long Short-Term Memory}.
Neural Computation, 9(8), 1735-1780.

\bibitem{rumelhart1986learning}
Rumelhart, D. E., Hinton, G. E., \& Williams, R. J. (1986).
\textit{Learning Representations by Back-Propagating Errors}.
Nature, 323(6088), 533-536.

\bibitem{nasapds2023}
NASA Planetary Data System. (2023).
\textit{Planetary Data System Imaging Atlas}.
Retrieved from https://pds-imaging.jpl.nasa.gov/

\bibitem{curiosity2023}
Mars Science Laboratory. (2023).
\textit{Curiosity Rover Mission Data}.
NASA Jet Propulsion Laboratory.

\bibitem{perseverance2023}
Mars 2020 Mission. (2023).
\textit{Perseverance Rover Mission Data}.
NASA Jet Propulsion Laboratory.

\bibitem{opencv2023}
Bradski, G. (2000).
\textit{The OpenCV Library}.
Dr. Dobb's Journal of Software Tools.

\bibitem{pytorch2023}
Paszke, A., Gross, S., Massa, F., Lerer, A., Bradbury, J., Chanan, G., \& Chintala, S. (2019).
\textit{PyTorch: An Imperative Style, High-Performance Deep Learning Library}.
In Advances in Neural Information Processing Systems (NeurIPS), 8024-8035.

\bibitem{streamlit2023}
Streamlit Inc. (2023).
\textit{Streamlit: The fastest way to build and share data apps}.
Retrieved from https://streamlit.io/

\end{thebibliography}

\end{document}


